\begin{theorem}\label{thm:ok-mvob-metrics}
Let \(L_i\) be big line bundles on \(X\), and let \(h_i\) be psh metrics
with \(\vol(\theta_i,\varphi_i)>0\), in the normalization of
\Cref{rem:ok-mvob-normalizations}.  Then
\[
 \vol(\ddc h_1,\dots,\ddc h_n)
 =n!\,\sup_\nu V\!\left(
    \Delta_\nu(L_1,h_1),\dots,\Delta_\nu(L_n,h_n)\right).
\tag{55.48}\label{eq:ok-mvob-metrics}
\]
\end{theorem}

\begin{proof}
For each \(i\), apply \Cref{lem:ok-mvob-i-model-approximation} and write
\(u_i=P_{\theta_i}[\varphi_i]_{\cI}\).  The replacement preserves every
body, and by definition \eqref{eq:ok-mvob-mixed-mass} the left side is the
mixed non-pluripolar mass of the \(u_i\).  Choose the approximants
\(u_{i,j}\) in \eqref{eq:ok-mvob-approximating}, using one common index for
the finitely many \(i\).  At every \(j\),
\Cref{prop:ok-mvob-analytic-case} gives
\[
 \int_X\bigwedge_i(\theta_i+\ddc u_{i,j})
 =n!\sup_\nu V(\Delta_\nu(\theta_1,u_{1,j}),\ldots,
                    \Delta_\nu(\theta_n,u_{n,j})).
\]
The mixed non-pluripolar masses on the left converge by
\cite[Theorem~6.2.1]{XiaBook}.  The right side converges by
\Cref{prop:ok-mvob-uniform-ds}; explicitly, for the corresponding functions
\(f_j,f\) of \(\nu\),
\[
 \left|\sup_\nu f_j(\nu)-\sup_\nu f(\nu)\right|
 \leq\sup_\nu|f_j(\nu)-f(\nu)|\longrightarrow0.
\]
Taking the limit proves \eqref{eq:ok-mvob-metrics}.
\end{proof}
